\chapter{Rachunek asymptotyczny}
\label{ma1:ch:asymptotics}

W rozdziale \ref{ma1:ch:taylor} stosowaliśmy już zapis małego $o$ do opisu reszty Peana. Rozwiniemy teraz ten język tak, aby porównywać funkcje zarówno w pobliżu punktu skończonego, jak i w nieskończoności. Asymptoty prostoliniowe omówiliśmy w rozdziale \ref{ma1:ch:function-variation}; rachunek asymptotyczny dotyczy również przybliżeń przez funkcje inne niż afiniczne.

\section{Symbole Landaua i równoważność asymptotyczna}
\begin{definicja}
\label{ma1:def:landau}
Niech $D\subseteq\R$ i niech $a$ będzie punktem skupienia $D$ lub niech $a=+\infty$ (odpowiednio $-\infty$), gdy $D$ jest nieograniczony z prawej (odpowiednio z lewej). Ustalmy kierunek graniczny $x\to a$, względnie do $D$; można także ograniczyć $D$ do jednej strony punktu. Wszystkie poniższe warunki dotyczą punktów $x\in D$, $x\ne a$, dostatecznie bliskich $a$, a dla granicy w nieskończoności dostatecznie dużych co do modułu w ustalonym kierunku. Niech $f,g:D\to\R$.
\begin{enumerate}
 \item Piszemy $f=O(g)$, jeżeli istnieje $C>0$ takie, że ostatecznie $|f(x)|\le C|g(x)|$.
 \item Piszemy $f=o(g)$, jeżeli dla każdego $\varepsilon>0$ ostatecznie $|f(x)|\le\varepsilon|g(x)|$.
 \item Piszemy $f=\Theta(g)$, jeżeli $f=O(g)$ oraz $g=O(f)$.
\end{enumerate}
Jeśli $g$ jest ostatecznie niezerowa, piszemy $f\sim g$, gdy $\lim_{x\to a}f(x)/g(x)=1$.
\end{definicja}
Słowo ,,ostatecznie'' oznacza istnienie odpowiedniego przebitego otoczenia punktu lub odpowiedniego ogona dziedziny. W definicji $o$ otoczenie może zależeć od $\varepsilon$. Definicje za pomocą nierówności mają sens także wtedy, gdy $g$ ma zera dowolnie blisko punktu granicznego. Na takich zerach odpowiednie oszacowanie wymusza $f(x)=0$. Dla ciągów używamy tych samych definicji, wybierając $D=\mathbb N$ oraz $a=+\infty$.

\begin{stwierdzenie}
\label{ma1:prop:landau-ratios}
Jeżeli $g$ jest ostatecznie niezerowa, to $f=O(g)$ jest równoważne ograniczoności $|f/g|$ w odpowiednim otoczeniu, a $f=o(g)$ jest równoważne $f/g\to0$. Ponadto
\[
 f\sim g \quad\Longleftrightarrow\quad f-g=o(g).
\]
W szczególności $f\sim g$ implikuje, że $f$ jest ostatecznie niezerowa i $f=\Theta(g)$.
\end{stwierdzenie}
\begin{proof}
Pierwsze dwie równoważności otrzymujemy, dzieląc nierówności przez dodatnią liczbę $|g(x)|$. Dla ostatniej korzystamy z tożsamości $(f-g)/g=f/g-1$. Jeśli $f/g\to1$, to ostatecznie $|f/g-1|<1/2$, skąd $1/2<f/g<3/2$. W szczególności $f\ne0$, $|f|\le(3/2)|g|$ i $|g|\le2|f|$.
\end{proof}
Zapisy $O(g)$ i $o(g)$ oznaczają klasy funkcji spełniających wskazane warunki. W zapisie $f=g+O(h)$ stwierdzamy, że $f-g=O(h)$; nie jest to zwykła równość dwóch konkretnych funkcji. Na przykład $x^2=O(x)$ przy $x\to0$, ale $x$ nie jest $O(x^2)$ w tej granicy. Symbole $O$ i $o$ opisują rząd oszacowania, a $\sim$ dodatkowo ustala współczynnik wiodący.

\section{Reguły rachunku}
\begin{twierdzenie}
\label{ma1:thm:landau-algebra}
Wszystkie relacje poniżej dotyczą tego samego kierunku granicznego. Zachodzą następujące reguły:
\begin{enumerate}
 \item Jeśli $f=O(g)$ i $h=O(g)$, to $\alpha f+\beta h=O(g)$ dla $\alpha,\beta\in\R$. To samo zachodzi dla małego $o$.
 \item Jeśli $f=O(g)$ i $h=O(k)$, to $fh=O(gk)$. Gdy jeden z dwóch warunków zastąpimy przez małe $o$, otrzymujemy $fh=o(gk)$.
 \item Jeśli $f=O(g)$ i $g=O(k)$, to $f=O(k)$. Gdy przynajmniej jeden warunek jest typu $o$, otrzymujemy $f=o(k)$.
 \item Jeśli $f=o(g)$, to $f=O(g)$.
\end{enumerate}
\end{twierdzenie}
\begin{proof}
Dla pierwszej reguły wybieramy wspólne otoczenie, w którym $|f|\le C|g|$ i $|h|\le B|g|$. Nierówność trójkąta daje
$|\alpha f+\beta h|\le(|\alpha|C+|\beta|B)|g|$; gdy współczynnik jest zerem, możemy zastąpić go dowolną dodatnią stałą. Dla małego $o$ przy zadanym $\varepsilon>0$ stosujemy do obu funkcji oszacowanie ze stałą $\varepsilon/(1+|\alpha|+|\beta|)$. Otrzymujemy $|\alpha f+\beta h|\le\varepsilon|g|$.

W drugiej regule mnożymy nierówności, otrzymując $|fh|\le CB|gk|$. Jeśli $f=o(g)$, dla zadanej $\varepsilon>0$ wybieramy oszacowanie $|f|\le(\varepsilon/B)|g|$ i $|h|\le B|k|$ z $B>0$. Ich iloczyn jest nie większy niż $\varepsilon|gk|$. Przypadek $h=o(k)$ wynika z przemienności. W trzeciej regule podstawiamy $|g|\le B|k|$ do $|f|\le C|g|$. Gdy pierwszy warunek jest typu $o$, zastępujemy $C$ przez $\varepsilon/B$, a gdy drugi jest typu $o$, zastępujemy $B$ przez $\varepsilon/C$. Ostatnia reguła wynika z wyboru $\varepsilon=1$.
\end{proof}

\begin{twierdzenie}
\label{ma1:thm:asymptotic-equivalence-algebra}
Na zbiorze funkcji ostatecznie niezerowych relacja $\sim$ jest relacją równoważności. Jeśli $f\sim g$ i $u\sim v$, to $fu\sim gv$ i $f/u\sim g/v$. Jeśli dodatkowo $g,v$ są ostatecznie tego samego znaku, to $f+u\sim g+v$.
\end{twierdzenie}
\begin{proof}
Zwrotność wynika z $f/f=1$. Symetria wynika z $g/f=1/(f/g)\to1$, a przechodniość z $f/h=(f/g)(g/h)\to1$. Dla iloczynu i ilorazu korzystamy odpowiednio z $(fu)/(gv)=(f/g)(u/v)$ oraz $(f/u)/(g/v)=(f/g)/(u/v)$. Wszystkie mianowniki są ostatecznie niezerowe. Dla sumy piszemy $f=g(1+r)$ i $u=v(1+s)$, gdzie $r,s\to0$. Wtedy
\[
 \left|\frac{f+u}{g+v}-1\right|
 \le\frac{|g||r|+|v||s|}{|g+v|}
 \le\max\{|r|,|s|\},
\]
gdyż $|g+v|=|g|+|v|>0$. Prawa strona dąży do zera, co dowodzi tezy.
\end{proof}
\begin{przyklad}
\label{ma1:ex:asymptotic-cancellation}
Nie wolno w ogólności zastępować składników sumy równoważnikami. Przy $x\to0^+$ weźmy $f=x+x^2$, $g=x$, $u=-x+2x^2$, $v=-x+x^2$. Mamy $f/g=1+x\to1$ oraz $u/v=(-1+2x)/(-1+x)\to1$, ale
\[
 \frac{f+u}{g+v}=\frac{3x^2}{x^2}=3.
\]
Wyrazy pierwszego rzędu uległy redukcji. Także różniczkowanie równoważników nie jest ogólnie uprawnione: dla $x>0$ niech $F(x)=x+x^2\sin(1/x)$, $G(x)=x$. Mamy $F/G=1+x\sin(1/x)\to1$, lecz
$F'(x)=1+2x\sin(1/x)-\cos(1/x)$, $G'(x)=1$. Dla $x_n=1/(2n\pi)$ mamy $F'(x_n)=0$, a dla $y_n=1/((2n+1)\pi)$ mamy $F'(y_n)=2$. Zatem $F'/G'$ nie ma granicy przy $x\to0^+$.
\end{przyklad}

\section{Porównywanie rzędów wzrostu}
\begin{twierdzenie}
\label{ma1:thm:asymptotic-growth}
Dla $p>0$, $q>0$, $c>0$ zachodzą relacje przy $x\to+\infty$:
\[
 (\log x)^p=o(x^q),\qquad x^p=o(e^{cx}).
\]
Ponadto, jeśli $\alpha<\beta$, to $x^\alpha=o(x^\beta)$ przy $x\to+\infty$ oraz $x^\beta=o(x^\alpha)$ przy $x\to0^+$.
\end{twierdzenie}
\begin{proof}
Wybieramy liczbę całkowitą $m>p$. Dla $t>0$ twierdzenie Taylora z resztą Lagrange'a \ref{ma1:thm:taylor-remainders} zastosowane do funkcji wykładniczej na $[0,t]$ daje $e^t\ge t^m/m!$, gdyż wszystkie składniki i reszta są dodatnie. Podstawiając $t=cx$, otrzymujemy
\[
 0\le \frac{x^p}{e^{cx}}\le\frac{m!}{c^m}x^{p-m}\longrightarrow0.
\]
Dla $x>1$ podstawiamy $t=\log x$ w już dowiedzionej relacji $t^p=o(e^{qt})$. Ponieważ $\log x\to+\infty$ i $e^{q\log x}=x^q$, otrzymujemy pierwszą tezę. Ostatnie dwie relacje wynikają z ilorazów $x^{\alpha-\beta}$ oraz $x^{\beta-\alpha}$ i granic dodatnich potęg.
\end{proof}
Dla ciągów otrzymujemy analogiczne relacje przez ograniczenie zmiennej do liczb naturalnych.
\begin{przyklad}
\label{ma1:ex:factorial-growth}
Dla każdej liczby $c>0$ mamy $c^n=o(n!)$. Istotnie, dla $a_n=c^n/n!$ zachodzi $a_{n+1}/a_n=c/(n+1)$. Wybieramy $N$ tak, aby dla $n\ge N$ iloraz był nie większy niż $1/2$. Indukcyjnie $0<a_{N+k}\le a_N2^{-k}$ dla $k\ge0$, więc $a_n\to0$.
\end{przyklad}

\begin{center}
\begin{tikzpicture}[x=.25cm,y=.16cm,>=Stealth]
 \draw[->] (0,0)--(22,0) node[right] {$x$};
 \draw[->] (0,0)--(0,34) node[above] {$y$};
 \draw[blue!65!black,thick,domain=1:20,samples=100] plot (\x,{ln(\x)});
 \draw[green!45!black,thick,domain=1:20,samples=100] plot (\x,{\x});
 \draw[orange!85!black,thick,domain=1:20,samples=100] plot (\x,{exp(\x/6)});
 \node[blue!65!black,above] at (17,2.83) {$\log x$};
 \node[green!45!black,above left] at (20,20) {$x$};
 \node[orange!85!black,above left] at (20,28.03) {$e^{x/6}$};
\end{tikzpicture}

\small Rysunek: przykładowe rzędy wzrostu. Wykres na skończonym odcinku ilustruje, ale nie dowodzi relacji $\log x=o(x)$ i $x=o(e^{x/6})$.
\end{center}

\section{Skale i rozwinięcia asymptotyczne}
\begin{definicja}
\label{ma1:def:asymptotic-scale}
W ustalonym kierunku granicznym ciąg funkcji $(\varphi_j)_{j\ge0}$, z których każda jest ostatecznie niezerowa, nazywamy skalą asymptotyczną, jeśli $\varphi_{j+1}=o(\varphi_j)$ dla każdego $j$. Mówimy, że $f$ ma rozwinięcie asymptotyczne w tej skali ze współczynnikami $(c_j)$, jeśli dla każdego całkowitego $N\ge0$
\[
 f(x)-\sum_{j=0}^N c_j\varphi_j(x)=o(\varphi_N(x)).
\]
Dla skończonego rozwinięcia wymagamy analogicznej własności do ustalonego rzędu $N$.
\end{definicja}
Zapis rozwinięcia nieskończonego nie oznacza zbieżności szeregu $\sum_{j\ge0}c_j\varphi_j(x)$ dla ustalonego $x$. Każda skończona suma daje osobne przybliżenie przy $x\to a$.

\begin{twierdzenie}
\label{ma1:thm:asymptotic-uniqueness}
Współczynniki rozwinięcia w ustalonej skali są jednoznacznie wyznaczone. Jeśli rozwinięcie istnieje do rzędu $N$, to dla $0\le k\le N$
\[
 c_k=\lim_{x\to a}
 \frac{f(x)-\sum_{j=0}^{k-1}c_j\varphi_j(x)}{\varphi_k(x)},
\]
gdzie dla $k=0$ suma jest pusta.
\end{twierdzenie}
\begin{proof}
Dla $j>k$ mamy $\varphi_j=o(\varphi_k)$ przez wielokrotne zastosowanie przechodniości z twierdzenia \ref{ma1:thm:landau-algebra}. Jeżeli reszta do rzędu $N$ jest $o(\varphi_N)$, to dla $k<N$ jest też $o(\varphi_k)$. Po odjęciu sumy do rzędu $k$ pozostają skończona suma składników $c_j\varphi_j$ dla $j>k$ oraz ta reszta, a zatem ich suma jest $o(\varphi_k)$. Otrzymujemy więc rozwinięcie także do każdego niższego rzędu. Dzieląc
$f-\sum_{j<k}c_j\varphi_j=c_k\varphi_k+o(\varphi_k)$ przez $\varphi_k$, otrzymujemy granicę z tezy. Dla $k=0$ wyznacza ona $c_0$ bez znajomości innych współczynników; kolejne granice wyznaczają współczynniki indukcyjnie, więc są jednoznaczne.
\end{proof}

\begin{przyklad}
\label{ma1:ex:asymptotic-geometric}
Dla $x\to0$ funkcja $1/(1-x)$ ma w skali $1,x,x^2,\ldots$ rozwinięcie o wszystkich współczynnikach równych $1$. Istotnie, dla $x\ne1$ tożsamość skończonej sumy geometrycznej daje
\[
 \frac1{1-x}=\sum_{j=0}^N x^j+\frac{x^{N+1}}{1-x}.
\]
Dla $|x|<1/2$ moduł reszty jest nie większy niż $2|x|^{N+1}$, więc jest ona $O(x^{N+1})$, a tym bardziej $o(x^N)$. Podstawiając $x=1/t$, otrzymujemy przy $t\to+\infty$
\[
 \frac1{t-1}=\sum_{j=1}^{N+1}\frac1{t^j}+O(t^{-N-2}).
\]
\end{przyklad}
Twierdzenie Taylora z resztą Peana \ref{ma1:thm:taylor-peano} jest przykładem skończonego rozwinięcia w skali $(x-a)^j$. Gdy funkcja jest gładka w otoczeniu $a$, otrzymujemy rozwinięcia do dowolnego rzędu. Nie wynika z tego odtwarzanie funkcji przez jej szereg Taylora: przykład \ref{ma1:ex:taylor-flat} ma wszystkie współczynniki zerowe, a funkcja nie jest zerowa poza punktem $a=0$.

\section{Obliczanie granic za pomocą rozwinięć}
\begin{przyklad}
\label{ma1:ex:asymptotic-cancel-taylor}
Obliczmy granicę
\[
 \lim_{x\to0}\frac{\sin x-x\cos x}{x^3}.
\]
Z twierdzenia \ref{ma1:thm:taylor-peano} i wzorów na pochodne trygonometryczne dostajemy
\[
 \sin x=x-\frac{x^3}{6}+o(x^3),\qquad
 \cos x=1-\frac{x^2}{2}+o(x^2).
\]
Mnożymy drugie rozwinięcie przez $x$, stosując regułę iloczynu z twierdzenia \ref{ma1:thm:landau-algebra}, i odejmujemy:
\[
 \sin x-x\cos x
 =\left(-\frac16+\frac12\right)x^3+o(x^3)
 =\frac{x^3}{3}+o(x^3).
\]
Po podzieleniu przez $x^3$ dostajemy granicę $1/3$. Samo zastąpienie $\sin x$ przez $x$ i $\cos x$ przez $1$ nie wystarczyłoby, gdyż wyrazy wiodące się redukują.
\end{przyklad}

\begin{przyklad}
\label{ma1:ex:asymptotic-reciprocal}
Jeżeli $\eta(x)\to0$, to ostatecznie $|\eta(x)|<1/2$. Tożsamość
\[
 \frac1{1+\eta}=1-\eta+\frac{\eta^2}{1+\eta}
\]
daje $1/(1+\eta)=1-\eta+O(\eta^2)$, gdyż $|1+\eta|\ge1/2$. Przy $x\to0$ załóżmy $\eta(x)=ax+O(x^2)$ dla stałej $a\in\R$. Wówczas $\eta=O(x)$: reszta $O(x^2)$ jest również $O(x)$ w otoczeniu zera. Z reguł rachunku $\eta^2=O(x^2)$ i
\[
 \frac1{1+\eta(x)}=1-ax+O(x^2).
\]
W szczególności z $\log(1+x)=x-x^2/2+o(x^2)$ wynika dla $x\ne0$ bliskiego zeru
\[
 \frac{x}{\log(1+x)}=1+\frac{x}{2}+o(x).
\]
Istotnie, $\log(1+x)/x=1-x/2+o(x)$, więc $\eta=-x/2+o(x)$ jest $O(x)$ i $O(\eta^2)=O(x^2)=o(x)$. Stosujemy poprzednią tożsamość. Dzielnik jest ostatecznie niezerowy, ponieważ $\log(1+x)/x\to1$.
\end{przyklad}
